% Supplementary material: the accounting boundary for environmental claims.

\section{What the Environmental Claim Covers}
\label{subsubsec:challenges_env_scope}
\label{supp:energy}

Environmental claims for small models are routinely made at a level of generality the accounting does not support. A statement that many small-model instances on consumer hardware use less total energy than one multi-GPU server, for example, cannot be sourced or derived: its truth depends on utilization on both sides, and a server at high batch utilization is considerably more energy-efficient per token than a consumer device performing occasional single-stream inference.

A complete comparison of a small-model deployment against a large-model one has to include the training energy of the small model; the training energy of any teacher used to distill it, which can dominate and which vendors do not report; the energy of the distillation, quantization, or pruning pipeline; the embodied carbon of the devices inference runs on, amortized over their service life and the share of it spent on inference; the marginal grid intensity where each side runs, which differs by more than a factor of ten between regions; and the energy of model updates, which for on-device models means redistributing weights to every device.

Almost none of these terms is reported. The available measurements are per-query and partial: a 2023 estimate put a frontier query at about ten times the energy of a web search \citep{devries2023energy}, while a 2025 production measurement covering accelerators, host systems, idle capacity, and data-centre overhead reports 0.24~Wh for the median text prompt of one large assistant \citep{elsworth2025google_energy}. Per-token compute scales with active parameters, so for a query both can answer a small model performs one to two orders of magnitude less compute than a frontier model; whether that becomes an energy saving depends on the utilization of the hardware each runs on (\evD). The defensible claim is therefore conditional: \emph{for a task a small model completes correctly, on hardware it keeps busy, serving it with the small model uses less inference energy}. Whether a shift to small models reduces total system energy is a question no published accounting answers, and the rebound effect, that cheaper inference induces more of it, cuts against it.
